\section{\texttt{TPARTMUL}}
\noindent\textbf{Descriptor profile.} Vec entry 39, tile catalog opcode \texttt{0x1044}. Descriptor bundle: \texttt{B.OP + B.ATTR + B.IOT}; WO/WM sites: 6; VEC beat-level sites: 6; ROM bytes: 48.\par\smallskip
\TileInstructionEncoding{figs/encoding/tile_instruction_entries/entry_39_pto_tpartmul_th0.pdf}{figs/encoding/tile_instruction_entries/entry_39_pto_tpartmul_th2.pdf}{0x1044}{39}{B.OP + B.ATTR + B.IOT}
\paragraph{Bundled descriptor (PTO-ISA header).}
\begin{description}[leftmargin=0.18\linewidth,style=nextline]
\item[Bundle] \texttt{B.OP + B.ATTR + B.IOT}. \texttt{B.OP.desc\_count} = 2; \texttt{B.OP.rom\_entry\_id} = 39.
\item[Assembly] 0: \texttt{B.OP} selects \texttt{TPARTMUL}, carries element format/variant, declares \texttt{desc\_count}, and names the ROM entry. 1: \texttt{B.ATTR} carries partial-valid-region policy. 2: \texttt{B.IOT} binds architectural tile operands. Across all \texttt{B.IOT} descriptors: tile inputs \texttt{src0}, \texttt{src1}, mask inputs none, and outputs \texttt{dst}.
\item[Tile operands] 2 tile input(s): \texttt{src0}, \texttt{src1}; 1 tile output(s): \texttt{dst}. Bound by 1 architectural \texttt{B.IOT} descriptor; each \texttt{B.IOT} supplies at most two source selectors plus one destination selector.
\item[Scalar GPR operands] No scalar GPR import/export descriptor is required.
\item[Metadata operands] \texttt{B.ATTR} carries partial-valid-region policy.
\end{description}
\par\smallskip
{\small
\paragraph{PTO ISA description.}
Partial elementwise multiply with implementation-defined handling of mismatched valid regions.
\paragraph{PTO ISA operation.}
For each element \texttt{(i, j)} in the destination valid region:

\[
\mathrm{dst}_{i,j} =
\begin{cases}
\mathrm{src0}_{i,j} \cdot \mathrm{src1}_{i,j} & \text{if both inputs are defined at } (i,j) \\
\mathrm{src0}_{i,j} & \text{if only src0 is defined at } (i,j) \\
\mathrm{src1}_{i,j} & \text{if only src1 is defined at } (i,j)
\end{cases}
\]
}\par\smallskip
\paragraph{Microcode site table.}
{\scriptsize
\begin{longtable}{@{}R{0.055\linewidth}L{0.23\linewidth}L{0.31\linewidth}L{0.22\linewidth}@{}}
\toprule
Seq & Micro instruction & WO[63:32] fields & WM[31:0] fields \\
\midrule
0 & \texttt{u\_vec.vlds} & \texttt{opc=0x17, x=0, fl=mem, seq=0, kind=b} & \texttt{entry=39, cnt=2, res=vec, var=0} \\
1 & \texttt{u\_vec.vlds} & \texttt{opc=0x17, x=0, fl=mem, seq=1, kind=b} & \texttt{entry=39, cnt=2, res=vec, var=0} \\
2 & \texttt{u\_vec.vmul} & \texttt{opc=0x1E, x=0, fl=alu, seq=2, kind=b} & \texttt{entry=39, cnt=2, res=vec, var=0} \\
3 & \texttt{u\_vec.vsts} & \texttt{opc=0x2A, x=0, fl=mem, seq=3, kind=b} & \texttt{entry=39, cnt=2, res=vec, var=0} \\
4 & \texttt{u\_vec.vlds} & \texttt{opc=0x17, x=0, fl=mem, seq=4, kind=b} & \texttt{entry=39, cnt=2, res=vec, var=0} \\
5 & \texttt{u\_vec.vsts} & \texttt{opc=0x2A, x=0, fl=mem, seq=5, kind=b} & \texttt{entry=39, cnt=2, res=vec, var=0} \\
\bottomrule
\end{longtable}
}
\paragraph{ROM body DSL.}
\begin{lstlisting}[style=ptocode]
def tpart_op_instr(dst, src0, src1, valid_rows, valid_cols):
    dtype = dst.element_type
    for row in range(0, valid_rows, 1):
        remained = valid_cols
        for col in range(0, valid_cols, pto.get_lanes(dtype)):
            mask, remained = pto.make_mask(dtype, remained)
            lhs = pto.vlds(src0[row, col:])
            rhs = pto.vlds(src1[row, col:])
            summed = pto.vmul(lhs, rhs, mask)
            pto.vsts(summed, dst[row, col:], mask)
    return None

def tpart_copy_instr(dst, src, valid_rows, valid_cols, start_row):
    dtype = dst.element_type
    for row in range(start_row, valid_rows, 1):
        remained = valid_cols
        for col in range(0, valid_cols, pto.get_lanes(dtype)):
            mask, remained = pto.make_mask(dtype, remained)
            val = pto.vlds(src[row, col:])
            pto.vsts(val, dst[row, col:], mask)
    return None

def tpart_op(dst, src0, src1,
             dst_valid_rows, dst_valid_cols,
             src1_valid_rows, src1_valid_cols):

    src1_eq_dst = (src1_valid_rows == dst_valid_rows and src1_valid_cols == dst_valid_cols)
    src1_row_lt_dst = (src1_valid_rows < dst_valid_rows and src1_valid_cols == dst_valid_cols)
    src1_col_lt_dst = (src1_valid_rows <= dst_valid_rows and src1_valid_cols < dst_valid_cols)

    if src1_eq_dst:
        tpart_op_instr(dst, src0, src1, dst_valid_rows, dst_valid_cols)
    elif src1_col_lt_dst:
        tpart_copy_instr(dst, src0, dst_valid_rows, dst_valid_cols, 0)
        if src1_valid_cols > 0:
            tpart_op_instr(dst, src0, src1, src1_valid_rows, src1_valid_cols)
    elif src1_row_lt_dst:
        if src1_valid_cols > 0:
            tpart_op_instr(dst, src0, src1, src1_valid_rows, src1_valid_cols)
        tpart_copy_instr(dst, src0, dst_valid_rows, dst_valid_cols, src1_valid_rows)

    return

def rom_tpartmul(src0, src1, dst):
    dst_valid_rows, dst_valid_cols = dst.valid_shape
    src0_valid_rows, src0_valid_cols = src0.valid_shape
    src1_valid_rows, src1_valid_cols = src1.valid_shape

    src0_eq_dst = (src0_valid_rows == dst_valid_rows and src0_valid_cols == dst_valid_cols)
    src1_eq_dst = (src1_valid_rows == dst_valid_rows and src1_valid_cols == dst_valid_cols)

    if src0_eq_dst or src1_eq_dst:
        if src0_eq_dst:
            tpart_op(dst, src0, src1, dst_valid_rows, dst_valid_cols, src1_valid_rows, src1_valid_cols)
        elif src1_eq_dst:
            tpart_op(dst, src1, src0, dst_valid_rows, dst_valid_cols, src0_valid_rows, src0_valid_cols)

    return
\end{lstlisting}
\Needspace{0.50\textheight}
